% Содержимое подраздела 1.1 - Представление молекулярных структур для задач машинного обучения

\section{Представление молекулярных структур для задач машинного обучения}
\label{sec:representation}

Ключевым шагом, позволившим применить методы обработки естественного языка (Natural Language Processing, NLP) к задачам медицинской химии, стала разработка линейных нотаций для представления молекулярных структур. Наиболее широкое распространение получила система \textbf{SMILES (Simplified Molecular-Input Line-Entry System)} --- формат, представляющий молекулу в виде строки символов ASCII. Данный подход был разработан с целью создания машиночитаемого, но при этом компактного и интуитивно понятного для человека представления~\cite{wigh2022review}.

Принцип кодирования структуры в SMILES основан на наборе простых правил:
\begin{itemize}
	\item \textbf{Атомы} представляются стандартными символами из периодической таблицы (например, \texttt{C} для углерода, \texttt{N} для азота). Атомы в ароматических системах обозначаются строчными буквами (например, бензол кодируется как \texttt{c1ccccc1}).
	\item \textbf{Связи} одинарные подразумеваются по умолчанию. Двойные и тройные связи обозначаются символами \texttt{=} и \texttt{\#} соответственно.
	\item \textbf{Ветвления} боковых цепей от основного скелета молекулы заключаются в круглые скобки \texttt{()}. Например, изобутан представляется строкой \texttt{CC(C)C}.
	\item \textbf{Циклы} кодируются путем разрыва одной из связей в цикле. Атомы, соединенные этой связью, помечаются одинаковыми цифровыми индексами. Например, циклогексан можно представить как \texttt{C1CCCCC1}, где цифра \texttt{1} указывает на замыкание цикла между первым и последним атомами углерода.
\end{itemize}

Представление SMILES обладает рядом преимуществ, главным из которых является его совместимость с моделями глубокого обучения. Рассматривая молекулу как <<предложение>> на <<химическом языке>>, исследователи смогли напрямую применить к задачам дизайна лекарств мощные архитектуры, такие как рекуррентные нейронные сети (RNN) и трансформеры. Кроме того, SMILES-строки компактны и более читаемы для человека по сравнению с другими форматами, такими как таблицы связности (connection tables), что упрощает работу с базами данных~\cite{wigh2022review}.

Тем не менее, у данного подхода есть и существенные недостатки. Основной проблемой является \textbf{неканоничность}: одна и та же молекула может быть представлена множеством синтаксически корректных SMILES-строк (например, \texttt{CCO} и \texttt{OCC} для этанола). Это усложняет сравнение структур и требует вычислительно затратных процедур канонизации. Другой важный недостаток заключается в том, что генеративные модели, обученные на SMILES, часто производят синтаксически некорректные строки (например, с непарными скобками). Это создает риск того, что модель в первую очередь изучает синтаксис нотации, а не фундаментальные химические закономерности. Наконец, в базовой форме SMILES не кодирует 3D-информацию о конформации молекулы, что является критически важным для многих задач.

Для решения этих проблем были разработаны альтернативные подходы. \textbf{InChI (International Chemical Identifier)} генерирует единственную, каноническую строку для любой молекулы, что решает проблему неоднозначности, но его сложная иерархическая структура делает его менее удобным для генеративных моделей. Более фундаментальной альтернативой являются \textbf{молекулярные графы}, где атомы представляются узлами, а связи -- ребрами. Такое представление является более естественным для молекул, лишено синтаксических артефактов и по построению не может сгенерировать химически невалидную структуру (например, с <<непарной связью>>). Однако работа с графовыми данными требует применения специализированных архитектур, таких как графовые нейронные сети.
